\section{Counting once, and checking the bars}
\label{sec:infaudit}

Two matters remain, and they are the chapter's bookkeeping and its
conscience: how to let the filter steer a real calibration without
counting any quote twice, and how to test --- against data --- that
the whole construction's stated uncertainty is true.

\subsection{One objective, one count}

\Cref{sec:inffilter} computed the update \emph{after} a data-only
fit: fit the quotes, read the handle, combine with the prediction.
When the filter is allowed to steer the fit itself, a tempting
implementation is fatal: build the posterior from today's quotes,
then feed that posterior back into the same fit as one more prior
row \emph{alongside those same quotes}.  The quotes then speak twice
--- once through the likelihood, once smuggled inside the
``prior'' --- and the published uncertainty is too small by exactly
the double count.  The correct implementation is a single objective:
today's quote rows, plus \emph{one} extra row per handle anchoring
the model's handle to the prediction $\hat O^{-}$ at weight
$1/\mathcal{V}^{-}$ --- the row form \cref{eq:infrow} with the
prediction as target.  That this is the \emph{same} answer as the
two-step update is a one-line calculation worth displaying:

\begin{proposition}[The joint fit is the filter]
\label{prop:infmap}
For a handle observed directly, the minimizer over $O$ of the joint
objective
\[
\frac{\bigl(O-\hat O_{\mathrm{obs}}\bigr)^2}{\mathcal{V}_{\mathrm{obs}}}
+\frac{\bigl(O-\hat O^{-}\bigr)^2}{\mathcal{V}^{-}}
\]
is the posterior $\hat O^{+}$ of \cref{eq:infupdate}, and the
objective's second derivative is twice the posterior precision
$1/\mathcal{V}^{+}$.
\end{proposition}

\begin{proof}
Set the derivative in $O$ to zero:
\[
\frac{2\bigl(O-\hat O_{\mathrm{obs}}\bigr)}{\mathcal{V}_{\mathrm{obs}}}
+\frac{2\bigl(O-\hat O^{-}\bigr)}{\mathcal{V}^{-}}=0
\;\Longrightarrow\;
O\Bigl(\tfrac{1}{\mathcal{V}_{\mathrm{obs}}}+\tfrac{1}{\mathcal{V}^{-}}\Bigr)
=\frac{\hat O_{\mathrm{obs}}}{\mathcal{V}_{\mathrm{obs}}}
+\frac{\hat O^{-}}{\mathcal{V}^{-}},
\]
the precision-weighted average --- which \cref{sec:inffilter}
already showed equals
$\hat O^{-}+\mathcal{K}(\hat O_{\mathrm{obs}}-\hat O^{-})$.  The
second derivative is
$2(1/\mathcal{V}_{\mathrm{obs}}+1/\mathcal{V}^{-})=2/\mathcal{V}^{+}$
by \cref{eq:infprecadd}.
\end{proof}

So the filter, done inside the objective, is one ungated row per
handle --- and the accounting rule follows: once the fit is
committed, its handles \emph{are} the posterior, and no second
update may be applied against the same quotes.  Note the word
\emph{ungated}: the prediction row carries its covariance weight
always, closed gates or no.  This is the chapter's two instruments
meeting in one objective, each keyed to its own symptom --- the
prior rows of \cref{sec:infgate} act where quotes are \emph{absent}
(their gate reads support), the filter row acts because quotes are
\emph{noisy} (its weight reads a variance) --- and they must not
overlap, because both pull toward the same transported yesterday.
Yesterday enters once.  When the filter runs, its prediction rows
\emph{are} yesterday's belly --- level, skew, curvature, transported
and priced --- so the shape baskets of
\cref{sec:infprice,sec:infcoords}, which carry that same yesterday
in other clothes, stand down; what persistence keeps is only what
the filter state does not transport: the deep-tail anchors beyond
the last basket leg, and --- for a node with no quotes at all ---
the baseline the next chapter builds.  Anchoring the same handle to
yesterday twice --- once through $\hat O^{-}$, once through a
persistence row --- would count yesterday twice, the mirror image of
the quote double count.

\subsection{The audit: were the error bars true?}

A filter can always make charts smoother.  Whether it is
\emph{right} is a different question, and it has a precise form:
the posterior comes with a stated variance, so the construction
claims that its errors, divided by their claimed scale, are of size
one.  Define the \emph{standardized error}
\begin{equation}
\mathcal{Z}
=\frac{\text{truth}-\hat O^{+}}{\sqrt{\mathcal{V}^{+}}},
\label{eq:infz}
\end{equation}
whose standard deviation should be $1$ if the error bars mean what
they say.  The test statistic is that standard deviation, taken over
a run of days: std$(\mathcal{Z})$ well above one convicts the filter
of overconfidence (bars too narrow); well below one, of timidity
(bars too wide).  The audit needs a truth to compare against, so it runs
wherever one exists --- held-out quotes on real history, or, as
here, a synthetic world where the truth is known exactly.

\begin{figure}[htbp]
\centering
\paperfig{\textwidth}{fig_flt_audit.pdf}
\caption{The audit on a seeded synthetic history
(\MacFiltAuditDays\ days; the true level diffuses at
\MacFiltAuditQTrueBp\ bp/$\sqrt{\text{day}}$ and jumps
$+\MacFiltAuditShockPts$ points twice; observations carry honestly
stated noise, wide on every seventh day).  (a)~A window around a
jump: the honest-budget filter (blue, with its $\pm1$-sd band)
follows in two days; the starved budget (violet) crawls.  (b)~The
audit curve: std of $\mathcal{Z}$ against the assumed walk scale ---
starved budgets sit far above the dotted target of $1$
(\MacFiltAuditZstdStarved\ at $10$~bp), the true scale reads
\MacFiltAuditZstdHonest.  (c)~Mean error against days since the
jump: starved \MacFiltAuditDayZeroStarved\ points on the day and
still \MacFiltAuditDayThreeStarved\ three days later; honest
\MacFiltAuditDayZeroHonest\ then pennies; the surprise widening cuts the
jump-day error to \MacFiltAuditDayZeroGate.}
\label{fig:infaudit}
\end{figure}

\Cref{fig:infaudit} runs the audit end to end on a seeded synthetic
history --- \MacFiltAuditDays\ days in which the true level
diffuses at exactly $\MacFiltAuditQTrueBp$ bp per square-root day,
jumps $+\MacFiltAuditShockPts$ volatility points on two days, and is
observed daily through honestly-stated noise
(Appendix~\ref{app:infprotocol} lists the generator).  The same
scalar filter \cref{eq:infupdate} is run with different assumed walk
scales $q_{\mathrm{walk}}$, and every claim of the last two sections
becomes measurable.

Panel~(a) is the texture.  Around the first jump, the
honest-budget filter (blue, its one-standard-deviation band shaded)
tracks the truth within about two days, its band wide enough to
contain its own errors.  The starved filter --- walk scale a third
of the truth --- is the violet curve: it enters the jump with a
narrow band, trusts its prediction too much, and crawls toward the
new level for a week.

Panel~(b) is the chapter's verdict in one curve: the standard
deviation of $\mathcal{Z}$, \cref{eq:infz}, against the assumed
walk scale.  At the true scale the curve sits near the dotted
target ---
std$(\mathcal{Z})=\MacFiltAuditZstdHonest$, the small excess being
the two jumps, which the plain clock term does not model; switch the
surprise widening on and it drops to \MacFiltAuditZstdGated.  To the
left the curve climbs steeply: at a third of the true scale the
standardized errors have standard deviation
\MacFiltAuditZstdStarved\ --- the filter's stated bars are three
times too narrow.  And panel~(c) shows the same starved filter's
\emph{speed}: \MacFiltAuditDayZeroStarved\ points of error on the
jump day, still \MacFiltAuditDayThreeStarved\ points three days
later, against the honest budget's
\MacFiltAuditDayZeroHonest-then-pennies, and
\MacFiltAuditDayZeroGate\ on the day with the surprise widening.  This
is the general lesson, and it is worth stating plainly because
intuition suggests the opposite: shrinking the prediction budget does
not make a filter cautious.  It makes it \emph{overconfident and
slow at once} --- sure of a surface it updates too slowly --- the
worst available combination, and the audit catches both symptoms in
one statistic.

\subsection{The contract}

The chapter closes as the others do: three columns of decreasing
strength --- proved, measured, chosen.

\begin{table}[htbp]
\centering
\small
\caption{The chapter's contract.}
\label{tab:infcontract}
\begin{tabular}{p{0.31\textwidth}p{0.31\textwidth}p{0.30\textwidth}}
\toprule
\textbf{Proved} & \textbf{Measured} & \textbf{Convention}\\
\midrule
The no-damp guarantee: a closed gate contributes nothing, for every
parameter vector (\cref{prop:infnodamp}). &
The flat valley: \MacPriorFlatValleyBp\ bp of rms movement across a
\MacPriorFlatPinSpanPts-point wing sweep on the thinned morning;
the full chain's V (\cref{fig:infflat}). &
The required precision, the gate exponent $\gamma$, and the prior
budget: set by judgment, disclosed, not derived.\\
\addlinespace
The harmonic law and the dead-leg rule; the RR/BF factor of
\MacPriorBasketFactor\ (\cref{prop:infharmonic}). &
Gates computed on the staged mornings: ATM
\MacPriorJumpGateAtm, moderate RR \MacPriorJumpGateRr\
(\cref{fig:infjump}). &
Support \cref{eq:infsupport} is a kernel proxy for identification;
it over-credits one bandwidth past the last quote
(\cref{sec:infblind}).\\
\addlinespace
Baskets annihilate the level direction; per-leg anchors penalize it
(\cref{prop:infnull}). &
The overnight experiment: basket wing error
\MacPriorJumpBasketErrDeep\ points, anchor
\MacPriorJumpAnchorErrDeep, data-only \MacPriorJumpDataErrDeep\
(\cref{fig:infjump}). &
Deep-tail anchors persist absolute prices beyond the last basket
leg: right when the level is stable, wrong across a jump ---
confined there by design.\\
\addlinespace
The update \cref{eq:infupdate} is the precision-weighted average;
the joint fit reproduces it exactly (\cref{prop:infmap}). &
The stale-strike morning: gains
\MacFiltCaseGainLevel/\MacFiltCaseGainSkew/\MacFiltCaseGainCurv\
computed from stated variances (\cref{fig:infupdate}). &
The state is three ATM handles, combined per handle: estimated
cross-handle correlations are deliberately not spent
(\cref{sec:inffilter}).\\
\addlinespace
The quadratic contract: observation variance is stated noise
squared, propagated (\cref{eq:infinfo}--\cref{eq:infdelta}). &
Slope \MacFiltCovarSlope\ in log--log; the contradiction sweep:
curvature gain \MacFiltCovarGainCurvClean$\to$%
\MacFiltCovarGainCurvKinked\ with the level barely moved
(\cref{fig:infcovar}). &
The prediction budget \cref{eq:infpred}: walk scale, market-clock
$\Delta t_{\mathrm{days}}$, transport and event widenings, the
surprise widening's threshold --- each an empirical commitment.\\
\addlinespace
--- (this row is empirical by design) &
The audit: std$(\mathcal{Z})$
\MacFiltAuditZstdStarved\ starved vs \MacFiltAuditZstdHonest\
honest (\MacFiltAuditZstdGated\ gated); jump-day errors
\MacFiltAuditDayZeroStarved/\MacFiltAuditDayZeroHonest/%
\MacFiltAuditDayZeroGate\ points (\cref{fig:infaudit}). &
std$(\mathcal{Z})=1$ is the acceptance test this book adopts for any
machinery that publishes uncertainty.\\
\bottomrule
\end{tabular}
\end{table}

Step back to the book's thesis.  \emph{Validity} in this chapter is
inherited, not new: every fit underneath carries the arbitrage
machinery of Part~I, and the filtered state is read off
arbitrage-checked smiles.  \emph{Information} is the chapter's
subject made explicit: for the first time the book has measured not
just what the data says but where it says \emph{nothing} --- the
flat valley --- and where it speaks with a noisy voice --- the
inflated budget.  \emph{Convention} is what fills the remainder,
and the chapter's discipline is that every convention entered
through a stated, auditable quantity: a gate with a dead zone, a
variance with a units contract, a walk scale that an audit can
convict.  Posterior equals data where the data speaks, and
yesterday --- transported, priced, and counted once --- where it
does not: the thesis of this book, made dynamic.

One boundary remains, and it is the book's last.  Everything here
is \emph{per node}: one smile, its own yesterday, its own morning.
But a desk's universe is not a list of unrelated smiles --- it is a
web of underliers and expiries whose laws move together, and a
morning that is sparse on one node is often rich on its neighbours.
What a lit node learns should reach the dark ones, at a price set
by how related they are --- the same precision language, now
flowing across a graph.  Building that graph, and completing a
whole universe of smiles from a sparse set of observed ones, is the
final chapter.
